\documentclass[aps,rmp,reprint,amsmath,amssymb,graphicx,longbibliography]{revtex4-1}

\usepackage{bm}
\usepackage{graphicx}
\usepackage{booktabs}
\usepackage{float}
\usepackage{algorithm}
\usepackage{algpseudocode}
\usepackage[utf8]{inputenc}
\usepackage{hyperref}
\hypersetup{breaklinks=true,colorlinks=true,linkcolor=blue,citecolor=blue,
            filecolor=magenta,urlcolor=blue}

\begin{document}

\title{Regression, resampling methods and gradient descent\\
       applied to Runge's function}

\author{Parisa Amin}
\thanks{Code and results: \url{https://github.com/USERNAME/FYS-STK4155}}
\affiliation{Department of Physics, University of Oslo, N-0316 Oslo, Norway}

\begin{abstract}
%% ~5 sentences, following the course guidelines:
%%   1. Short introduction to the topic and why it matters
%%   2. The challenge / unresolved issue you address
%%   3. What you did to address it
%%   4. Main results -- WITH NUMBERS
%%   5. The implications
\end{abstract}

\maketitle

%% ===================================================================
\section{Introduction}
%% - Motivate the reader: why polynomial regression, why Runge's function,
%%   why regularisation and resampling matter.
%% - What I have done.
%% - Close with a paragraph on how the report is organised
%%   ("In Sec.~\ref{sec:theory} we ... in Sec.~\ref{sec:results} we ...").

%% ===================================================================
\section{Theory and methods}\label{sec:theory}

\subsection{Data: Runge's function}
%% f(x) = 1/(1+25x^2) on [-1,1], design matrix, noise model y = f(x) + eps.

\subsection{Ordinary least squares}
%% Cost function, normal equations, the closed-form solution via SVD/pinv.
%% The analytical gradient (week 35 exercises).

\subsection{Ridge regression}
%% The l2 penalty, closed form, and the SVD picture: shrinkage of the
%% singular-value modes (Sec. 3.8 of the lecture notes).

\subsection{Lasso regression}
%% The l1 penalty; why there is no closed form; the subgradient at zero.

\subsection{Scaling and the train--test split}
%% Why we centre/standardise, and what it does to the intercept.

\subsection{The bias--variance decomposition}
%% DERIVATION -- required by part c). E[(y-ytilde)^2] = Bias^2 + Var + sigma^2.
%% State clearly over which randomness the expectations are taken, and discuss
%% what replacing f by y does to the measured bias.

\subsection{Resampling: the bootstrap and \textit{k}-fold cross-validation}

\subsection{Gradient descent methods}
%% Plain GD and the learning-rate bound eta < 2/lambda_max(H).
%% Momentum, AdaGrad, RMSprop, Adam. Stochastic gradient descent.

\subsection{Automatic differentiation}
%% Why AD is neither symbolic nor finite-difference differentiation, and its
%% cost relative to one evaluation of the cost function.

%% ===================================================================
\section{Implementation and validation}\label{sec:code}
%% Structure of the code (modules, functions). Then the BENCHMARKS -- the
%% course guidelines call this "extremely important":
%%   (i)   gradient descent reproduces the closed-form OLS/Ridge solution
%%   (ii)  our Ridge matches scikit-learn (mind the alpha convention)
%%   (iii) the analytical gradient matches the JAX gradient to machine precision

%% ===================================================================
\section{Results and discussion}\label{sec:results}

\subsection{OLS on Runge's function}
\subsection{Ridge regression and the role of $\lambda$}
\subsection{Bias--variance tradeoff from the bootstrap}
\subsection{Cross-validation, and comparison with the bootstrap}
\subsection{Gradient descent: convergence and the learning rate}
\subsection{Momentum and adaptive learning rates}
\subsection{Lasso}
\subsection{Stochastic gradient descent}
\subsection{Final model selection: OLS vs.\ Ridge vs.\ Lasso}

%% ===================================================================
\section{Conclusions and perspectives}
%% Main findings and interpretations. Pros and cons of each method.
%% Possible improvements and future work.

%% ===================================================================
\section*{Declaration on the use of large language models}
%% MANDATORY -- 5 points, and omitting it costs 5 more. Follow
%% https://github.com/EducationalMaterialUiO/MachineLearningUiO/blob/main/LLM_Usage_Declaration_Guidelines.md

%% ===================================================================
\appendix
\section{Additional derivations and figures}

\bibliography{References}

\end{document}
